\section{Related Works}
\label{sec:ch-6-related_work}

As surveyed in~\Cref{sec:background-eval-summaries}, standard automatic summarization metrics show low correlation with human judgment~\cite{liu2008correlation,fabbri2021summeval}, motivating ongoing work on metrics that better align with human preferences. However, there is little work in reproducing studies on evaluation metrics, despite our field's reliance on their outcomes. We include a number of traditional metrics in both our reproduction study and as baselines for our LLM-as-judge experiments. 

 Recent work has shown great success using LLMs-as-judges, showing high correlation with human judgments across a variety of tasks and domains~\cite{zheng2023judging}. Researchers have improved LLM-as-Judge performance through methods such as evaluation-specific finetuning~\cite{kim2023prometheus} or better prompting~\cite{liu2023geval}. Although LLMs-as-Judges perform best when provided with a gold reference~\cite{Deutsch2022LimitationsReferenceFree}, there is great interest in the reference-free setting of evaluation metrics as gold references are often unavailable or prohibitively expensive to collect~\cite{liu2023geval,vasilyev2020blanc}. We propose a new setting, reference-adjusted evaluation, in which relevant but non-exact references are available.